\documentclass[12pt,a4paper]{article}

% ==========================================
% IDIOMA Y TIPOGRAFÍA
% ==========================================

\usepackage[spanish]{babel}
\usepackage{fontspec}
\setmainfont{Arial}
\setmonofont{Menlo}[Scale=0.85]

% ==========================================
% MÁRGENES Y ESPACIADO
% ==========================================

\usepackage[
    top=2.5cm,
    bottom=2.5cm,
    left=2.5cm,
    right=2.5cm
]{geometry}

\usepackage{setspace}
\setstretch{1.15}

\setlength{\parindent}{0pt}
\setlength{\parskip}{6pt}

% ==========================================
% MATEMÁTICAS
% ==========================================

\usepackage{amsmath}
\usepackage{amssymb}
\usepackage{mathtools}

% ==========================================
% IMÁGENES Y GRÁFICOS
% ==========================================

\usepackage{graphicx}
\usepackage{float}
\usepackage{tikz}
\usepackage{pgfplots}

\pgfplotsset{
    compat=1.18,
    every axis/.append style={
        /pgf/number format/assume math mode=true,
        /pgf/number format/use comma,
        /pgf/number format/1000 sep={}
    }
}
\graphicspath{{images/}}

\usetikzlibrary{
    arrows.meta,
    positioning,
    shapes.geometric,
    fit,
    backgrounds,
    calc
}

% Colores de cada mecanismo, iguales en todo el documento
\definecolor{rest}{HTML}{1565C0}
\definecolor{kafka}{HTML}{6A1B9A}
\definecolor{grpc}{HTML}{EF6C00}
\definecolor{ws}{HTML}{2E7D32}
\definecolor{falla}{HTML}{C62828}
\definecolor{gris}{HTML}{616161}
\definecolor{fondo}{HTML}{F5F5F5}

\tikzset{
    flecha/.style={-{Stealth[length=2.2mm]}, thick},
    caja/.style={draw, rounded corners=2pt, fill=white, align=center,
        font=\small, inner sep=4pt, minimum height=0.9cm},
    servicio/.style={caja, minimum width=2.6cm},
    actor/.style={caja, fill=fondo, minimum width=2.6cm},
    topico/.style={draw, cylinder, shape border rotate=90, aspect=0.18,
        fill=kafka!10, draw=kafka, align=center, font=\small,
        minimum height=1.3cm, minimum width=2.6cm},
    estado/.style={draw, rounded corners=6pt, fill=white, align=center,
        font=\footnotesize\ttfamily, inner sep=4pt, minimum height=0.75cm},
    nota/.style={font=\footnotesize, text=gris, align=center}
}

% ==========================================
% TABLAS
% ==========================================

\usepackage{booktabs}
\usepackage{array}
\usepackage{longtable}
\usepackage{tabularx}
\usepackage{caption}

\renewcommand{\arraystretch}{1.2}
\newcolumntype{Y}{>{\raggedright\arraybackslash}X}

\captionsetup{
    font=small,
    labelfont=bf
}
\captionsetup[table]{position=bottom}

% Texto en monoespaciado para nombres de código
\newcommand{\cod}[1]{\texttt{#1}}

% ==========================================
% BIBLIOGRAFÍA - APA 7
% ==========================================

\usepackage{csquotes}

\usepackage[
    backend=biber,
    style=apa
]{biblatex}

\addbibresource{referencias.bib}
\appto{\bibsetup}{\raggedright}

% ==========================================
% CONFIGURACIÓN GENERAL
% ==========================================

\usepackage[hidelinks]{hyperref}

\pagestyle{empty}

% ==========================================
% INICIO DEL DOCUMENTO
% ==========================================

\begin{document}

% ==========================================
% PORTADA
% ==========================================

\hypersetup{pageanchor=false}
\begin{titlepage}
\centering
\thispagestyle{empty}

% ==========================================
% LOGO INSTITUCIONAL
% ==========================================

\vspace*{0.5cm}

\includegraphics[width=10cm]{images/logo_ups.png}

\vspace{1.2cm}

% ==========================================
% INFORMACIÓN INSTITUCIONAL
% ==========================================

{\Large\bfseries
UNIVERSIDAD POLITÉCNICA SALESIANA
\par}

\vspace{0.4cm}

{\large
SEDE CUENCA
\par}

\vspace{0.6cm}

{\large
INGENIERÍA EN CIENCIAS DE LA COMPUTACIÓN
\par}

% ==========================================
% TÍTULO DEL TRABAJO
% ==========================================

\vfill

\rule{0.75\textwidth}{0.5pt}

\vspace{0.6cm}

{\LARGE\bfseries
CAMPUSBITE: MECANISMOS DE COMUNICACIÓN EN UN SISTEMA DISTRIBUIDO DE PEDIDOS
\par}

\vspace{0.6cm}

\rule{0.75\textwidth}{0.5pt}

\vfill

% ==========================================
% INFORMACIÓN ACADÉMICA
% ==========================================

\begin{minipage}{0.85\textwidth}
\raggedright
\large

\textbf{Asignatura:}
Sistemas Distribuidos

\vspace{0.5cm}

\textbf{Docente:}
Ing. Cristian Timbi

\vspace{0.5cm}

\textbf{Estudiantes:}
David Larriva, Carolina Fortmann, Domenica Uyunkar y Cinthya Ramón

\end{minipage}

% ==========================================
% LUGAR Y FECHA
% ==========================================

\vfill

{\normalsize
Cuenca, Ecuador
\par}

\vspace{0.2cm}

{\normalsize
\today
\par}

\vspace{0.3cm}

\end{titlepage}
\hypersetup{pageanchor=true}

% ==========================================
% ÍNDICE
% ==========================================

\pagestyle{plain}
\setcounter{page}{1}

{\setstretch{1.0}\tableofcontents}
\clearpage

% ==========================================
% CONTENIDO PRINCIPAL
% ==========================================

\section{Introducción}

En las horas pico los estudiantes hacen fila para pedir, para pagar y para esperar su comida, y no saben cuándo estará listo su pedido. CampusBite es un sistema distribuido que resuelve este problema para la cafetería del campus. El estudiante pide desde la web o el celular, cada pedido llega a cocina, caja y facturación, el estudiante recibe un aviso cuando su pedido está listo y la administración puede ver las estadísticas de ventas.

La práctica pide usar REST para la API de pedidos y otros mecanismos de comunicación del Capítulo 2, con al menos uno asíncrono y todos dentro del mismo flujo \parencite{timbi2026comunicacion}. Como el grupo tiene cuatro integrantes usamos cuatro mecanismos, uno por persona. La Figura~\ref{fig:resumen} muestra cómo un mismo pedido los recorre.

\begin{figure}[H]
\centering
\begin{tikzpicture}[node distance=0.6cm and 1.1cm]
    \node[actor] (est) {Estudiante};
    \node[servicio, right=1.6cm of est] (api) {API de Pedidos};
    \node[topico, right=1.6cm of api] (k) {Kafka};
    \node[servicio, above right=0.2cm and 1.4cm of k] (coc) {Cocina};
    \node[servicio, right=1.4cm of k] (caja) {Caja};
    \node[servicio, below right=0.2cm and 1.4cm of k] (fac) {Facturación};
    \node[servicio, below=1.4cm of api] (ws) {Notificaciones};

    \draw[flecha, rest] (est) -- node[above, font=\scriptsize]{REST} (api);
    \draw[flecha, kafka] (api) -- node[above, font=\scriptsize]{publica} (k);
    \draw[flecha, kafka] (k) -- (coc);
    \draw[flecha, kafka] (k) -- (caja);
    \draw[flecha, kafka] (k) -- (fac);
    \draw[flecha, grpc, very thick] (fac.east) -| node[pos=0.25, below, font=\scriptsize]{gRPC} ([xshift=0.5cm]caja.east) -- (caja.east);
    \draw[flecha, kafka] (k.south) |- (ws.east);
    \draw[flecha, ws, dashed] (ws.west) -| node[pos=0.25, below, font=\scriptsize]{WebSocket} (est.south);
\end{tikzpicture}
\caption{Recorrido de un pedido por los cuatro mecanismos. Cada color es un mecanismo y se usa igual en todo el documento: azul REST, morado Kafka, naranja gRPC y verde WebSocket.}
\label{fig:resumen}
\end{figure}

\subsection{Objetivo y alcance}

El objetivo es diseñar e implementar en Python un sistema donde cada mecanismo se use donde tiene sentido técnico, y comprobar con pruebas reales cómo responde ante tres fallas que pide la consigna: un pedido enviado dos veces, un servicio interno que no responde y un cliente que se desconecta.

El sistema tiene seis servicios y guarda los datos en memoria, algo que la consigna permite. Kafka corre en Docker y es lo único que persiste. No hay autenticación ni sistemas externos, porque el pago se hace en el mostrador y el cajero solo lo registra.

\begin{table}[H]
\centering
\small
\begin{tabularx}{\textwidth}{@{}l l Y@{}}
\toprule
\textbf{Integrante} & \textbf{Mecanismo} & \textbf{Servicios a cargo} \\
\midrule
David Larriva & REST & API de Pedidos, contrato, schema y arquitectura \\
Carolina Fortmann & Kafka & Broker, capa común de Kafka, Cocina y Estadísticas \\
Domenica Uyunkar & gRPC & Caja (REST y servidor gRPC) y Facturación \\
Cinthya Ramón & WebSocket & Notificaciones, cliente web e integración \\
\bottomrule
\end{tabularx}
\caption{Reparto del trabajo. Cada integrante es dueño de un mecanismo y de sus servicios.}
\label{tab:equipo}
\end{table}

\section{Arquitectura}

Antes de programar definimos la arquitectura con una versión corta de la metodología de la Práctica 01. Tiene tres partes, el diagrama de contexto, los escenarios de calidad priorizados y el mapa de microservicios.

\subsection{Diagrama de contexto (C4 nivel 1)}

El diagrama de contexto muestra quién usa CampusBite y con qué se comunica, sin entrar en cómo está construido por dentro. Hay tres actores y ningún sistema externo.

\begin{figure}[H]
\centering
\begin{tikzpicture}
    \node[draw=rest!70!black, fill=rest!85!black, text=white, rounded corners=3pt,
          align=center, minimum width=6.8cm, minimum height=2cm, font=\small] (cb) at (0,0)
          {\textbf{CampusBite}\\[2pt]\footnotesize Recibe pedidos, los reparte a cocina, caja y\\[-2pt]\footnotesize facturación, avisa al estudiante y calcula ventas};
    \node[draw, fill=rest!15, rounded corners=3pt, align=center, font=\small,
          minimum width=3.6cm, minimum height=1.3cm] (est) at (-5,3.4)
          {\textbf{Estudiante}\\\footnotesize pide desde web o celular};
    \node[draw, fill=rest!15, rounded corners=3pt, align=center, font=\small,
          minimum width=3.6cm, minimum height=1.3cm] (caj) at (5,3.4)
          {\textbf{Cajero}\\\footnotesize cobra y entrega};
    \node[draw, fill=rest!15, rounded corners=3pt, align=center, font=\small,
          minimum width=3.6cm, minimum height=1.3cm] (adm) at (0,-3.4)
          {\textbf{Administración}\\\footnotesize revisa las ventas};

    \draw[flecha, rest] ([xshift=-10mm]est.south) |- node[pos=0.25, left, font=\scriptsize, align=right]{crea y consulta\\pedidos\\REST} ([yshift=4mm]cb.west);
    \draw[flecha, ws, dashed] ([xshift=-18mm]cb.north) |- node[pos=0.25, right, font=\scriptsize, align=left]{avisa cambios\\WebSocket} (est.east);
    \draw[flecha, rest] (caj.south) |- node[pos=0.25, right, font=\scriptsize, align=left]{registra pago\\y entrega\\REST} (cb.east);
    \draw[flecha, rest] (adm.north) -- node[right, font=\scriptsize]{consulta estadísticas, REST} (cb.south);
\end{tikzpicture}
\caption{Diagrama de contexto C4 nivel 1. Kafka y gRPC no aparecen porque son internos y se ven en el mapa de microservicios.}
\label{fig:c4}
\end{figure}

\subsection{Escenarios de calidad}

Definimos seis escenarios con el formato de la Práctica 01, que tiene fuente, estímulo, artefacto, entorno, respuesta y medida. Para priorizarlos calificamos la importancia para la cafetería y la dificultad técnica, como en un árbol de utilidad. Los que tienen las dos altas son los que justifican los mecanismos elegidos (Figura~\ref{fig:prioridad}).

\begin{figure}[H]
\centering
\begin{tikzpicture}[esc/.style={circle, draw=rest, fill=rest!20, minimum size=0.62cm, inner sep=0pt, font=\small\bfseries}]
\begin{axis}[
    width=9cm, height=7cm,
    xmin=0.5, xmax=3.5, ymin=0.5, ymax=3.5,
    xtick={1,2,3}, xticklabels={Baja, Media, Alta},
    ytick={1,2,3}, yticklabels={Baja, Media, Alta},
    xlabel={Dificultad técnica}, ylabel={Importancia para la cafetería},
    tick label style={font=\small}, label style={font=\small},
    grid=major, grid style={gris!25},
]
    \fill[falla!8] (axis cs:2.5,2.5) rectangle (axis cs:3.5,3.5);
    \node[font=\scriptsize, text=falla] at (axis cs:3,3.38) {prioridad máxima};
    \node[esc] at (axis cs:2.82,3) {1};
    \node[esc] at (axis cs:3.18,3) {2};
    \node[esc] at (axis cs:2,3) {3};
    \node[esc] at (axis cs:1.82,2) {4};
    \node[esc] at (axis cs:2.18,2) {5};
    \node[esc] at (axis cs:1,2) {6};
\end{axis}
\end{tikzpicture}
\caption{Priorización de los seis escenarios, numerados como en el Cuadro~\ref{tab:escenarios}. Arriba a la derecha están los que, si fallan, hacen que el sistema mienta o cobre dos veces.}
\label{fig:prioridad}
\end{figure}

\begin{table}[H]
\centering
\small
\begin{tabularx}{\textwidth}{@{}c l Y l l@{}}
\toprule
\textbf{\#} & \textbf{Atributo} & \textbf{Estímulo y medida} & \textbf{Mecanismo} & \textbf{Prueba} \\
\midrule
1 & Disponibilidad & Kafka caído al crear un pedido. Responder \cod{503} en $\le$ 5 s y nunca \cod{201} sin evento & REST + Kafka & REST, parte 2 \\
2 & Confiabilidad & El cliente reintenta el mismo pedido. Un solo evento por clave & REST & REST, parte 1 \\
3 & Resiliencia & Caja no responde a Facturación. Cada llamada $\le$ 2 s, máximo 6 intentos & gRPC & gRPC \\
4 & Rendimiento & Pedidos con consumidores detenidos. La API responde igual y no se pierde ningún evento & Kafka & Kafka \\
5 & Recuperación & El estudiante pierde el WebSocket. Reintento cada 3 s y estado correcto al volver & WebSocket + REST & WebSocket \\
6 & Interoperabilidad & Evento que no cumple el contrato. 0 estados dañados & Kafka + contrato & Automáticas \\
\bottomrule
\end{tabularx}
\caption{Escenarios de calidad priorizados, con su medida y la prueba que los verifica.}
\label{tab:escenarios}
\end{table}

\subsection{Mapa de microservicios}

CampusBite sigue una arquitectura de microservicios \parencite{richardson2018microservices} con seis servicios. Cada uno guarda sus datos en memoria y es dueño de una parte del estado del pedido. Ningún servicio cambia los datos de otro, solo publica eventos o llama a su interfaz. La Figura~\ref{fig:mapa} muestra cada comunicación con su mecanismo y si es síncrona o asíncrona.

\begin{figure}[H]
\centering
\begin{tikzpicture}[scale=0.95, transform shape]
    % Actores
    \node[actor, minimum width=2.4cm] (cli) at (0,0) {Estudiante};
    \node[actor, minimum width=2.4cm] (caj) at (0,-2.2) {Cajero};
    \node[actor, minimum width=2.4cm] (adm) at (0,-4.4) {Administración};

    % Servicios del lado del cliente
    \node[servicio] (ws) at (4.2,0) {Notificaciones\\\scriptsize WebSocket :8765};
    \node[servicio] (api) at (4.2,-2.2) {API de Pedidos\\\scriptsize REST :8000};
    \node[servicio] (est) at (4.2,-4.4) {Estadísticas\\\scriptsize REST :8002};

    % Kafka en el centro
    \node[draw=kafka, fill=kafka!8, rounded corners=4pt, minimum width=1.6cm,
          minimum height=5.6cm, align=center, font=\small, text=kafka] (bus) at (7.6,-2.2)
          {\textbf{Kafka}\\\scriptsize :9092\\[8pt]\scriptsize pedido.\\[-3pt]\scriptsize creado\\[6pt]\scriptsize pedido.\\[-3pt]\scriptsize estado};

    % Servicios internos
    \node[servicio] (fac) at (11,0) {Facturación\\\scriptsize cliente gRPC};
    \node[servicio] (caja) at (11,-2.2) {Caja\\\scriptsize REST :8001 · gRPC :50051};
    \node[servicio] (coc) at (11,-4.4) {Cocina};

    % REST
    \draw[flecha, rest] (cli) -- (api);
    \draw[flecha, rest] ([yshift=2mm]caj.east) -- ([yshift=2mm]caj.east -| api.west);
    \draw[flecha, rest] (adm) -- (est);
    \draw[flecha, rest] ([yshift=-2mm]caj.east) -- (1.75,-2.4) -- (1.75,-5.75) -- (13.2,-5.75) |- (caja.east);
    % WebSocket
    \draw[{Stealth[length=2.2mm]}-{Stealth[length=2.2mm]}, thick, ws, dashed] (cli) -- (ws);
    % gRPC
    \draw[flecha, grpc, very thick] (fac) -- node[right, font=\scriptsize]{gRPC} (caja);
    % Kafka
    \draw[flecha, kafka] ([yshift=1.5mm]api.east) -- ([yshift=1.5mm]api.east -| bus.west);
    \draw[flecha, kafka] ([yshift=-1.5mm]api.east -| bus.west) -- ([yshift=-1.5mm]api.east);
    \draw[flecha, kafka] (ws.east -| bus.west) -- (ws.east);
    \draw[flecha, kafka] (est.east -| bus.west) -- (est.east);
    \draw[flecha, kafka] (fac.west -| bus.east) -- (fac.west);
    \draw[flecha, kafka] ([yshift=1.5mm]caja.west -| bus.east) -- ([yshift=1.5mm]caja.west);
    \draw[flecha, kafka] ([yshift=-1.5mm]caja.west) -- ([yshift=-1.5mm]caja.west -| bus.east);
    \draw[flecha, kafka] ([yshift=1.5mm]coc.west -| bus.east) -- ([yshift=1.5mm]coc.west);
    \draw[flecha, kafka] ([yshift=-1.5mm]coc.west) -- ([yshift=-1.5mm]coc.west -| bus.east);

    % Leyenda
    \matrix[column sep=5pt, font=\scriptsize, anchor=north] at (6.6,-6.2) {
        \draw[flecha, rest] (0,0) -- (0.7,0); & \node[anchor=west]{REST, síncrono}; &
        \draw[flecha, kafka] (0,0) -- (0.7,0); & \node[anchor=west]{Kafka, asíncrono}; &
        \draw[flecha, grpc, very thick] (0,0) -- (0.7,0); & \node[anchor=west]{gRPC, síncrono}; &
        \draw[thick, ws, dashed] (0,0) -- (0.7,0); & \node[anchor=west]{WebSocket, asíncrono}; \\
    };
\end{tikzpicture}
\caption{Mapa de microservicios. Kafka conecta a todos los servicios por eventos. REST y gRPC son las únicas llamadas directas, y el WebSocket queda abierto entre Notificaciones y el estudiante.}
\label{fig:mapa}
\end{figure}

\begin{table}[H]
\centering
\small
\begin{tabularx}{\textwidth}{@{}l c Y l@{}}
\toprule
\textbf{Servicio} & \textbf{Puerto} & \textbf{Responsabilidad} & \textbf{Estado que cambia} \\
\midrule
API de Pedidos & 8000 & Valida, aplica la idempotencia, calcula el total y publica \cod{pedido.creado} & \cod{RECIBIDO}, \cod{ENTREGADO} \\
Cocina & no tiene & Prepara el pedido sin esperar el pago & \cod{EN\_PREPARACION}, \cod{LISTO} \\
Caja & 8001 y 50051 & Registra el cobro, recibe el pago y responde \cod{ConfirmarPago} & \cod{PENDIENTE}, \cod{PAGADO} \\
Facturación & no tiene & Con \cod{PAGO/PAGADO} confirma con Caja y emite una factura & \cod{FACTURADO} \\
Notificaciones & 8765 & Reenvía a cada cliente los eventos de su pedido & ninguno \\
Estadísticas & 8002 & Cuenta pedidos, pagos, entregas y ventas & ninguno \\
\bottomrule
\end{tabularx}
\caption{Servicios, puertos y dueño de cada estado. Cada estado tiene un solo dueño, así se sabe quién hizo cada cambio.}
\label{tab:servicios}
\end{table}

\subsection{Flujo de un pedido}

La Figura~\ref{fig:secuencia} sigue un pedido de principio a fin. Las flechas continuas son llamadas síncronas, donde quien llama espera la respuesta, y las de Kafka son eventos que cada servicio procesa a su ritmo.

\begin{figure}[H]
\centering
\begin{tikzpicture}[
    yscale=0.66,
    msg/.style={flecha, font=\scriptsize},
    lbl/.style={font=\scriptsize, inner sep=1.5pt}
]
    \foreach \x/\n in {0/Cliente, 2.3/API, 4.6/Kafka, 6.9/Cocina, 9.2/Caja, 11.5/Facturación, 13.8/Notificac.} {
        \node[caja, minimum width=1.9cm, minimum height=0.6cm, font=\scriptsize] at (\x,0.7) {\n};
        \draw[gris!50, dashed] (\x,0) -- (\x,-16);
    }
    \draw[msg, rest] (0,-0.6) -- node[lbl, above]{1. POST /pedidos} (2.3,-0.6);
    \draw[msg, kafka] (2.3,-1.5) -- node[lbl, above]{2. pedido.creado} (4.6,-1.5);
    \draw[msg, rest, dashed] (2.3,-2.3) -- node[lbl, above]{201} (0,-2.3);
    \draw[msg, ws, dashed] (0,-3.2) -- node[lbl, above, pos=0.5]{3. suscribir por WebSocket y después GET} (13.8,-3.2);
    \draw[msg, kafka] (4.6,-4.1) -- node[lbl, above]{4. lee} (6.9,-4.1);
    \draw[msg, kafka] (4.6,-4.9) -- node[lbl, above, pos=0.75]{4. lee} (9.2,-4.9);
    \draw[msg, kafka] (9.2,-5.8) -- node[lbl, above, pos=0.3]{PENDIENTE} (4.6,-5.8);
    \draw[msg, kafka] (6.9,-6.7) -- node[lbl, above]{5. EN\_PREPARACION} (4.6,-6.7);
    \draw[msg, kafka] (4.6,-7.6) -- node[lbl, above, pos=0.8]{reenvío} (13.8,-7.6);
    \draw[msg, ws] (13.8,-8.5) -- node[lbl, above, pos=0.5]{push al estudiante} (0,-8.5);
    \draw[msg, rest] (0,-9.5) -- node[lbl, above, pos=0.55]{6. el cajero cobra: POST /caja/pagos} (9.2,-9.5);
    \draw[msg, kafka] (9.2,-10.4) -- node[lbl, above, pos=0.4]{PAGADO} (4.6,-10.4);
    \draw[msg, kafka] (4.6,-11.3) -- node[lbl, above, pos=0.8]{7. lee PAGADO} (11.5,-11.3);
    \draw[msg, grpc, very thick] (11.5,-12.2) -- node[lbl, above]{ConfirmarPago} (9.2,-12.2);
    \draw[msg, grpc, dashed] (9.2,-13.0) -- node[lbl, above]{PAGADO} (11.5,-13.0);
    \draw[msg, kafka] (11.5,-13.9) -- node[lbl, above, pos=0.2]{FACTURADO} (4.6,-13.9);
    \draw[msg, kafka] (6.9,-14.8) -- node[lbl, above]{8. LISTO} (4.6,-14.8);
    \draw[msg, rest] (0,-15.7) -- node[lbl, above]{9. entregar} (2.3,-15.7);
    \draw[msg, kafka] (2.3,-15.7) -- node[lbl, above]{ENTREGADO} (4.6,-15.7);
\end{tikzpicture}
\caption{Secuencia de un pedido. Cada cambio de estado pasa por \cod{pedido.estado} y llega por WebSocket al estudiante, como en el paso 5. Para no saturar el diagrama solo se dibuja ese reenvío una vez.}
\label{fig:secuencia}
\end{figure}

\subsection{Estados del pedido}

El estado tiene cuatro partes independientes y cada una tiene un solo servicio con autoridad para cambiarla. Un estado nunca retrocede, y si llega un evento con una combinación que no está en el contrato se ignora.

\begin{figure}[H]
\centering
\begin{tikzpicture}
    \node[font=\small\bfseries, text=gris, anchor=west] at (-0.2,0.8) {Preparación};
    \node[estado] (r) at (0.8,0) {RECIBIDO};
    \node[estado] (p) at (4.2,0) {EN\_PREPARACION};
    \node[estado] (l) at (7.6,0) {LISTO};
    \node[estado, fill=ws!15] (en) at (11.6,0) {ENTREGADO};
    \draw[flecha] (r) -- node[above, font=\scriptsize]{Cocina} (p);
    \draw[flecha] (p) -- node[above, font=\scriptsize]{Cocina} (l);
    \draw[flecha] (l) -- node[above, font=\scriptsize]{API} (en);

    \node[font=\small\bfseries, text=gris, anchor=west] at (-0.2,-1.8) {Pago y facturación};
    \node[estado] (pe) at (0.8,-2.6) {PENDIENTE};
    \node[estado] (pa) at (4.2,-2.6) {PAGADO};
    \node[estado] (fa) at (9.2,-2.0) {FACTURADO};
    \node[estado] (fp) at (9.2,-3.3) {FACTURACION\_PENDIENTE};
    \draw[flecha] (pe) -- node[above, font=\scriptsize]{Caja} (pa);
    \draw[flecha] (pa.east) -- node[above, font=\scriptsize, pos=0.35, sloped]{confirma} (fa.west);
    \draw[flecha] (pa.east) -- node[below, font=\scriptsize, pos=0.35, sloped]{6 fallos} (fp.west);

    \draw[flecha, dashed, gris] (pa.north) -- (4.2,-1.0) -| node[pos=0.25, above, font=\scriptsize]{la entrega también exige PAGADO} (en.south);
    \node[font=\scriptsize, text=gris, anchor=west] at (2.2,-1.8) {};
\end{tikzpicture}
\caption{Estados del pedido. La preparación la cambia Cocina (la API solo pone \cod{RECIBIDO}), el pago lo cambia Caja, la facturación la cambia Facturación y la entrega la API, que solo la acepta con el pedido \cod{LISTO} y \cod{PAGADO}. \cod{FACTURACION\_PENDIENTE} es final en esta versión.}
\label{fig:estados}
\end{figure}

\section{Mecanismos de comunicación y decisiones}

Cada mecanismo tiene su ADR en el repositorio. Aquí resumimos qué decidimos, por qué y qué alternativas descartamos. El criterio general viene de la guía de selección del Capítulo 2 \parencite{timbi2026comunicacion}: lo síncrono va donde alguien necesita una respuesta para seguir, y lo asíncrono donde basta con que la información llegue.

\begin{table}[H]
\centering
\small
\begin{tabularx}{\textwidth}{@{}l l Y Y@{}}
\toprule
\textbf{Mecanismo} & \textbf{Tipo} & \textbf{Dónde se usa} & \textbf{Por qué ahí} \\
\midrule
\textcolor{rest}{\textbf{REST}} & Síncrono & Cliente y cajero con la API, cajero con Caja, administración con Estadísticas & El estudiante necesita su \cod{idPedido} enseguida y HTTP funciona en cualquier navegador \\
\textcolor{kafka}{\textbf{Kafka}} & Asíncrono, persistente & Entre todos los servicios, con \cod{pedido.creado} y \cod{pedido.estado} & La API no espera a nadie y un servicio caído recupera los eventos al volver \\
\textcolor{grpc}{\textbf{gRPC}} & Síncrono con deadline & Facturación pregunta a Caja con \cod{ConfirmarPago} & Hace falta una respuesta confirmada antes de emitir la factura \\
\textcolor{ws}{\textbf{WebSocket}} & Asíncrono, server push & Notificaciones avisa al estudiante & El servidor avisa apenas cambia el estado, sin que el cliente consulte \\
\bottomrule
\end{tabularx}
\caption{Resumen de los cuatro mecanismos (ADR-005).}
\label{tab:mecanismos}
\end{table}

\subsection{REST: API de Pedidos (ADR-001)}

La API sigue el estilo REST \parencite{fielding2000architectural}, está hecha con FastAPI y tiene cuatro rutas. El cliente solo manda el producto y la cantidad, y el total lo calcula el servidor con el menú, en centavos enteros para no tener errores de redondeo.

\begin{table}[H]
\centering
\small
\begin{tabularx}{\textwidth}{@{}l l Y@{}}
\toprule
\textbf{Ruta} & \textbf{Éxito} & \textbf{Errores principales} \\
\midrule
\cod{GET /menu} & \cod{200} & ninguno \\
\cod{POST /pedidos} & \cod{201} nuevo, \cod{200} reintento & \cod{400} datos o falta la clave, \cod{409} en curso o clave reutilizada, \cod{503} Kafka no confirmó \\
\cod{GET /pedidos/\{id\}} & \cod{200} & \cod{404} no existe \\
\cod{PATCH /pedidos/\{id\}/entregar} & \cod{200}, también si se repite & \cod{409} no está \cod{LISTO} o no está \cod{PAGADO} \\
\bottomrule
\end{tabularx}
\caption{Rutas de la API de Pedidos. Todos los errores traen los campos \cod{error} y \cod{mensaje}.}
\label{tab:rutas}
\end{table}

La parte más delicada es no duplicar pedidos. El cliente genera una \cod{Idempotency-Key} por pedido y la API la guarda en un registro con un lock. La Figura~\ref{fig:clave} muestra los estados de la clave.

\begin{figure}[H]
\centering
\begin{tikzpicture}[node distance=1.2cm and 3.1cm]
    \node[estado, minimum width=2.6cm] (nueva) {clave nueva};
    \node[estado, minimum width=2.6cm, right=of nueva, fill=grpc!12] (proc) {EN\_PROCESO};
    \node[estado, minimum width=2.6cm, right=of proc, fill=ws!15] (comp) {COMPLETADA};
    \node[estado, minimum width=2.6cm, below=1.6cm of proc, fill=falla!10] (lib) {LIBERADA};

    \draw[flecha] (nueva) -- node[above, font=\scriptsize]{se reserva con lock} (proc);
    \draw[flecha] (proc) -- node[above, font=\scriptsize]{Kafka confirma} node[below, font=\scriptsize]{responde \cod{201}} (comp);
    \draw[flecha] ([xshift=-4mm]proc.south) -- node[left, font=\scriptsize, align=right, pos=0.4]{Kafka no confirma\\en 5 s, \cod{503}} ([xshift=-4mm]lib.north);
    \draw[flecha] ([xshift=4mm]lib.north) -- node[right, font=\scriptsize, align=left]{reintento con\\el mismo \cod{idPedido}} ([xshift=4mm]proc.south);

    \node[nota, above=0.35cm of proc] {otra solicitud igual recibe\\\cod{409 SOLICITUD\_EN\_CURSO}};
    \node[nota, below=0.35cm of comp] {un reintento recibe \cod{200}\\con la respuesta guardada};
    \node[nota, left=1.2cm of lib, align=right] {si la clave llega con otro cuerpo\\(hash SHA-256 distinto) recibe\\\cod{409 CLAVE\_REUTILIZADA}};
\end{tikzpicture}
\caption{Estados de la clave de idempotencia. \cod{LIBERADA} conserva el \cod{idPedido} porque un \cod{503} no prueba que Kafka no haya escrito el mensaje. Si sí lo escribió, los consumidores descartan el duplicado por \cod{idPedido}.}
\label{fig:clave}
\end{figure}

Descartamos gRPC para el cliente porque los navegadores necesitarían gRPC-Web y un proxy, crear pedidos por WebSocket porque habría que inventar un protocolo de petición y respuesta, y aceptar con \cod{202} para publicar después porque, con datos en memoria, ese pedido se perdería al reiniciar.

\subsection{Kafka: comunicación entre servicios (ADR-002)}

Kafka corre en Docker en modo KRaft y guarda los eventos en un log persistente \parencite{kafka2024documentation}. Hay dos topics de tres particiones, creados a propósito porque la creación automática usaría una sola, y un tercer topic \cod{pedido.dlq} para los mensajes que fallan después de reintentar.

\begin{figure}[H]
\centering
\begin{tikzpicture}[
    part/.style={draw=kafka, fill=kafka!6, minimum height=0.55cm, minimum width=6.2cm, anchor=west},
    ev/.style={draw=kafka, fill=kafka!35, minimum size=0.42cm, inner sep=0pt, font=\tiny}
]
    \node[font=\small\bfseries, text=kafka, anchor=west] at (0,1.1) {Topic \cod{pedido.estado}};
    \foreach \i/\y in {0/0.3, 1/-0.45, 2/-1.2} {
        \node[part] (p\i) at (0,\y) {};
        \node[font=\scriptsize, text=kafka, anchor=east] at (6.1,\y) {P\i};
    }
    \foreach \k/\t in {0/EP, 1/PE, 2/LI, 3/PA, 4/FA, 5/EN} {
        \node[ev] at (0.4+\k*0.55,0.3) {\t};
    }
    \node[ev, fill=gris!20, draw=gris] at (0.4,-0.45) {};
    \node[ev, fill=gris!20, draw=gris] at (0.95,-0.45) {};
    \node[ev, fill=gris!20, draw=gris] at (0.4,-1.2) {};
    \draw[flecha, gris] (0.2,-1.75) -- node[below, font=\scriptsize]{los offsets crecen hacia la derecha} (6.0,-1.75);

    \node[caja, font=\scriptsize, anchor=east, align=left] (prod) at (-0.9,-0.45)
        {Cocina, Caja,\\Facturación y API\\publican con clave\\\cod{idPedido}};
    \draw[flecha, kafka] (prod.east) -- (-0.45,-0.45) |- (p0.west);

    \foreach \g/\y in {api/0.9, notificaciones/0.15, facturacion/-0.6, estadisticas/-1.35} {
        \node[caja, font=\scriptsize, minimum width=2.6cm, minimum height=0.55cm, anchor=west] (g\g) at (7.2,\y) {grupo \cod{\g}};
        \draw[flecha, kafka] (6.25,-0.45) -- (g\g.west);
    }
\end{tikzpicture}
\caption{Particiones (P0, P1 y P2) y consumer groups. Todos los eventos de un pedido llevan la misma clave y caen en la misma partición, donde Kafka garantiza el orden. Cada grupo lee el topic completo, así que el mismo evento les llega a los cuatro. EP, PE, LI, PA, FA y EN son los seis estados de un pedido.}
\label{fig:particiones}
\end{figure}

\begin{table}[H]
\centering
\small
\begin{tabularx}{\textwidth}{@{}l c l Y@{}}
\toprule
\textbf{Topic} & \textbf{Particiones} & \textbf{Publican} & \textbf{Consumer groups} \\
\midrule
\cod{pedido.creado} & 3 & API & \cod{cocina}, \cod{caja}, \cod{estadisticas} \\
\cod{pedido.estado} & 3 & API, Cocina, Caja, Facturación & \cod{api}, \cod{notificaciones}, \cod{facturacion}, \cod{estadisticas} \\
\cod{pedido.dlq} & 1 & cualquier consumidor & mensajes que fallaron tras los reintentos \\
\bottomrule
\end{tabularx}
\caption{Topics del sistema.}
\label{tab:topics}
\end{table}

El productor usa \cod{acks='all'} y espera la confirmación hasta 5 s. El consumidor hace el commit del offset recién después de procesar el mensaje, así que si se cae a la mitad lo vuelve a leer. Eso hace que Kafka entregue al menos una vez, y por eso cada consumidor descarta los repetidos, \cod{pedido.creado} por \cod{idPedido} y \cod{pedido.estado} por \cod{idEvento}. Descartamos REST entre servicios porque la API tendría que esperar a cada uno, y Sockets TCP porque habría que construir a mano el protocolo y los mensajes se perderían con el receptor caído.

\subsection{gRPC: Facturación consulta a Caja (ADR-003)}

Antes de emitir una factura, Facturación necesita saber por la fuente autorizada si el pago está confirmado. Para eso usa una llamada gRPC unaria definida en \cod{caja.proto} con Protocol Buffers, lo que da un contrato tipado y código generado para cliente y servidor \parencite{grpc2024documentation}.

\begin{figure}[H]
\centering
\begin{tikzpicture}[
    msg/.style={draw, rounded corners=2pt, fill=grpc!8, draw=grpc, align=left, font=\scriptsize\ttfamily, inner sep=5pt}
]
    \node[msg] (req) {ConfirmarPagoRequest\\\quad id\_pedido};
    \node[caja, right=2.2cm of req, minimum width=3.4cm] (srv) {\textbf{Caja}\\\scriptsize rpc ConfirmarPago};
    \node[msg, right=1.0cm of srv] (res) {ConfirmarPagoResponse\\\quad estado: PAGADO | PENDIENTE\\\quad\quad\quad | NO\_ENCONTRADO\\\quad id\_confirmacion, monto\\\quad repetida};
    \draw[flecha, grpc, very thick] (req) -- node[above, font=\scriptsize]{deadline 2 s} (srv);
    \draw[flecha, grpc, very thick] (srv) -- (res);
\end{tikzpicture}
\caption{Contrato de \cod{ConfirmarPago}. El campo \cod{repetida} indica que esa confirmación ya se había entregado, lo que da la semántica a lo sumo una vez.}
\label{fig:proto}
\end{figure}

Si Caja no responde, el deadline corta la llamada a los 2 s con \cod{DEADLINE\_EXCEEDED} y Facturación reintenta cada 10 s, hasta seis veces. La Figura~\ref{fig:reintentos} muestra el peor caso, cuando todas las llamadas agotan el deadline.

\begin{figure}[H]
\centering
\begin{tikzpicture}[x=0.2cm, y=1cm]
    \draw[flecha, gris] (0,0) -- (66,0) node[right, font=\scriptsize]{s};
    \foreach \t in {0,12,24,36,48,60} {
        \fill[grpc] (\t,0.15) rectangle (\t+2,0.55);
        \node[font=\tiny, above] at (\t+1,0.55) {2 s};
        \draw[gris] (\t,0.05) -- (\t,-0.05) node[below, font=\scriptsize]{\t};
    }
    \foreach \t in {2,14,26,38,50} {
        \draw[gris, dashed] (\t,0.35) -- (\t+10,0.35);
        \node[font=\tiny, text=gris, above] at (\t+5,0.35) {espera 10 s};
    }
    \foreach \i/\t in {1/0,2/12,3/24,4/36,5/48,6/60} {
        \node[font=\scriptsize, text=grpc] at (\t+1,-0.7) {\i};
    }
    \node[font=\scriptsize, text=grpc, anchor=east] at (-1,-0.7) {intento};
    \draw[flecha, falla] (62,0.35) -- (62,1.25) node[above, font=\scriptsize, text=falla, align=center]{\cod{FACTURACION\_PENDIENTE}\\a los 62 s};
\end{tikzpicture}
\caption{Política de reintentos de Facturación en el peor caso: 6 llamadas de 2 s y 5 esperas de 10 s suman 62 s. Como el consumidor es secuencial, otros pedidos pueden esperar ese tiempo.}
\label{fig:reintentos}
\end{figure}

Una prueba nos mostró un problema de diseño. Al principio Facturación consumía \cod{pedido.creado} y empezaba a preguntar apenas se creaba el pedido. Si el estudiante pagaba después del último intento, el pedido se quedaba sin factura aunque estuviera pagado. En la versión 1.3 del contrato Facturación se activa con el evento \cod{PAGO/PAGADO}, que Kafka guarda hasta que alguien lo procese (Figura~\ref{fig:v13}).

\begin{figure}[H]
\centering
\begin{tikzpicture}[x=0.11cm, y=1cm]
    % Antes
    \node[font=\small\bfseries, anchor=west] at (0,1.55) {Antes: se activaba con \cod{pedido.creado}};
    \draw[gris] (0,0) -- (120,0);
    \fill[grpc!70] (0,-0.2) rectangle (62,0.2);
    \node[font=\scriptsize, text=white] at (31,0) {6 intentos, Caja responde PENDIENTE};
    \node[draw=falla, fill=falla!10, font=\scriptsize, rounded corners=2pt] (fp) at (78,0.62) {FACTURACION\_PENDIENTE};
    \draw[flecha, falla] (62,0.2) -- (fp.south west);
    \draw[flecha, rest] (95,0.9) node[above, font=\scriptsize]{el estudiante paga} -- (95,0.25);
    \node[font=\scriptsize, text=falla] at (100,-0.5) {nadie factura};
    \node[font=\scriptsize, text=gris] at (0,-0.45) {crea};

    % Después
    \node[font=\small\bfseries, anchor=west] at (0,-1.4) {Ahora: se activa con \cod{PAGO/PAGADO}};
    \draw[gris] (0,-2.4) -- (120,-2.4);
    \node[font=\scriptsize, text=gris] at (0,-2.85) {crea};
    \draw[flecha, rest] (95,-1.5) node[above, font=\scriptsize]{} -- (95,-2.15);
    \node[font=\scriptsize, anchor=east, text=rest] at (94,-1.75) {paga, Caja publica PAGADO};
    \fill[grpc!70] (95,-2.6) rectangle (99,-2.2);
    \node[draw=ws, fill=ws!15, font=\scriptsize, rounded corners=2pt, anchor=west] at (101,-2.4) {FACTURADO};
    \node[font=\scriptsize, text=grpc] at (97,-2.95) {gRPC};
\end{tikzpicture}
\caption{Por qué cambiamos el disparador de Facturación. El evento solo avisa que hubo un pago, y gRPC sigue confirmándolo con Caja antes de facturar.}
\label{fig:v13}
\end{figure}

Descartamos usar Kafka también para esta consulta, porque habría que publicar la pregunta y esperar la respuesta en otro topic, y gRPC entre todos los servicios, porque volvería síncrono lo que hoy es asíncrono. RPC/RMI clásico serviría, pero sin contrato tipado ni deadline por llamada, y RMI exige Java en los dos lados.

\subsection{WebSocket: notificaciones al estudiante (ADR-004)}

El servidor de notificaciones lee \cod{pedido.estado} con su propio consumer group y le reenvía cada evento solo a los clientes suscritos a ese pedido. La conexión empieza como HTTP con \cod{Upgrade: websocket}, el servidor responde \cod{101 Switching Protocols} y desde ahí queda abierta en las dos direcciones \parencite{mdn2024websockets}.

El WebSocket es transitorio. Si el estudiante no está conectado el aviso se descarta, pero el estado no se pierde porque lo tiene la API. Por eso el cliente sigue un orden fijo, primero se suscribe y después consulta con \cod{GET}. La Figura~\ref{fig:orden} muestra qué pasaría con el orden contrario.

\begin{figure}[H]
\centering
\begin{tikzpicture}[x=1cm, y=1cm]
    \foreach \y/\titulo/\a/\b/\col in {0/{Orden incorrecto: primero GET}/GET/suscribir/falla, -2.3/{Orden correcto: primero suscribir}/suscribir/GET/ws} {
        \node[font=\small\bfseries, anchor=west, text=\col] at (0,\y+0.75) {\titulo};
        \draw[flecha, gris] (0,\y) -- (11,\y) node[right, font=\scriptsize]{t};
        \draw[thick] (1.5,\y-0.12) -- (1.5,\y+0.12) node[above, font=\scriptsize]{\cod{\a}};
        \draw[thick] (7.5,\y-0.12) -- (7.5,\y+0.12) node[above, font=\scriptsize]{\cod{\b}};
        \fill[kafka] (4.5,\y) circle (3pt) node[below=2pt, font=\scriptsize, text=kafka]{llega \cod{LISTO}};
    }
    \node[font=\scriptsize, text=falla, align=left, anchor=west] at (8.3,-0.45) {el GET ya pasó y aún no\\había suscripción: se pierde};
    \node[font=\scriptsize, text=ws, align=left, anchor=west] at (8.3,-2.75) {ya estaba suscrito: llega por\\push, y lo anterior viene en el GET};
\end{tikzpicture}
\caption{Por qué el cliente se suscribe antes de consultar. Así todo cambio llega por una de las dos vías. Si la conexión se cae, el cliente reintenta cada 3 s y repite los dos pasos.}
\label{fig:orden}
\end{figure}

Descartamos el polling con \cod{GET} cada pocos segundos porque genera tráfico aunque el pedido no cambie y avisa con retraso, Server-Sent Events porque es de un solo sentido, y gRPC streaming porque no funciona directo en el navegador.

\subsection{Semánticas de entrega}

Las tres operaciones con efectos, que son crear el pedido, pagar y entregar, se pueden repetir por un reintento. El sistema combina dos semánticas del Capítulo 2 para que el efecto ocurra una sola vez \parencite{tanenbaum2017distributed,coulouris2012distributed}.

\begin{figure}[H]
\centering
\begin{tikzpicture}[node distance=0.3cm and 1.3cm]
    \node[caja, fill=rest!10, minimum width=3.3cm] (c) {\textbf{Cliente}\\\scriptsize reintenta si no hay respuesta};
    \node[caja, fill=grpc!10, right=of c, minimum width=3.3cm, minimum height=1.6cm] (f) {\textbf{Servidor}\\\scriptsize descarta duplicados\\\scriptsize por clave, \cod{idPedido}\\\scriptsize o \cod{idEvento}};
    \node[caja, fill=ws!15, right=of f, minimum width=2.6cm] (r) {\textbf{Un solo efecto}};
    \draw[flecha] ([yshift=4mm]c.east) -- ([yshift=4mm]f.west);
    \draw[flecha] (c.east) -- (f.west);
    \draw[flecha] ([yshift=-4mm]c.east) -- ([yshift=-4mm]f.west);
    \draw[flecha] (f) -- (r);
    \node[nota, below=0.25cm of c] {al menos una vez};
    \node[nota, below=0.25cm of f] {a lo sumo una vez};
    \node[nota, below=0.25cm of r] {resultado};
\end{tikzpicture}
\caption{El cliente repite, el servidor filtra y el resultado es un solo pedido, un solo pago y una sola entrega.}
\label{fig:semanticas}
\end{figure}

\section{Pruebas de comportamiento}

Hicimos una prueba por mecanismo ante una falla o una situación límite, más una prueba de integración con los seis servicios. Todas se ejecutaron con Kafka real en Docker y los resultados son los que obtuvimos. El Cuadro~\ref{tab:pruebas} resume lo que se probó y la sección de cada prueba muestra los datos.

\begin{table}[H]
\centering
\small
\begin{tabularx}{\textwidth}{@{}l c Y Y@{}}
\toprule
\textbf{Prueba} & \textbf{Escenario} & \textbf{Falla provocada} & \textbf{Resultado} \\
\midrule
REST, parte 1 & 2 & El mismo pedido enviado varias veces y a la vez & 7 solicitudes, 2 pedidos y 2 eventos \\
REST, parte 2 & 1 & Kafka detenido al crear un pedido & \cod{503} en 5,00 s y reintento con el mismo \cod{idPedido} \\
Kafka & 4 & Cocina detenida mientras llega un pedido & Cocina lo procesó al volver, sin perderlo \\
gRPC & 3 & Caja sin responder & Llamada cortada en 2,02 s con \cod{DEADLINE\_EXCEEDED} \\
WebSocket & 5 & El estudiante se desconecta antes de \cod{LISTO} & Recuperó \cod{LISTO} por REST y recibió \cod{ENTREGADO} por push \\
Integración & 6 & Pago y entrega repetidos en el flujo completo & 7 eventos en orden y un solo pago, factura y entrega \\
\bottomrule
\end{tabularx}
\caption{Resumen de las pruebas. Todas cumplieron la medida de su escenario.}
\label{tab:pruebas}
\end{table}

\subsection{REST: pedido duplicado y Kafka caído}

\textbf{Parte 1.} Un script envió el mismo pedido dos veces, cinco solicitudes simultáneas con la misma clave, una clave reutilizada con otro pedido y una solicitud sin clave. Al final leyó \cod{pedido.creado} desde el inicio para contar los eventos, que es lo que ven cocina, caja y estadísticas.

\begin{figure}[H]
\centering
\begin{tikzpicture}
\begin{axis}[
    width=6.4cm, height=5.6cm,
    ybar, bar width=0.6cm,
    ymin=0, ymax=8.5,
    symbolic x coords={Solicitudes, Pedidos, Eventos},
    xtick=data, enlarge x limits=0.25,
    xticklabel style={font=\scriptsize},
    nodes near coords, nodes near coords style={font=\small\bfseries},
    tick label style={font=\small}, ytick=\empty, axis x line*=bottom, axis y line=none,
    title={\small Lo enviado frente a lo creado},
]
    \addplot[fill=rest!70, draw=rest] coordinates {(Solicitudes,7) (Pedidos,2) (Eventos,2)};
\end{axis}
\end{tikzpicture}
\hfill
\begin{tikzpicture}
\begin{axis}[
    width=8.4cm, height=5.6cm,
    xbar stacked, bar width=0.55cm,
    xmin=0, xmax=5,
    symbolic y coords={Ejecución anterior, Ejecuciones 1 a 3},
    ytick=data, enlarge y limits=0.5,
    tick label style={font=\scriptsize}, xtick={0,1,...,5},
    xlabel={\scriptsize solicitudes simultáneas},
    legend style={font=\scriptsize, at={(0.5,1.03)}, anchor=south, legend columns=3, draw=none},
    title={\small 5 envíos a la vez con la misma clave}, title style={yshift=4mm},
]
    \addplot[fill=ws!60, draw=ws] coordinates {(1,Ejecuciones 1 a 3) (1,Ejecución anterior)};
    \addplot[fill=rest!40, draw=rest] coordinates {(0,Ejecuciones 1 a 3) (3,Ejecución anterior)};
    \addplot[fill=grpc!50, draw=grpc] coordinates {(4,Ejecuciones 1 a 3) (1,Ejecución anterior)};
    \legend{\cod{201} creado, \cod{200} reintento, \cod{409} en curso}
\end{axis}
\end{tikzpicture}
\caption{Resultados de la parte 1. A la izquierda, siete solicitudes con dos claves dejaron dos pedidos y dos eventos. A la derecha, el reparto de respuestas cambia según cuándo llega cada solicitud, pero siempre hay un solo \cod{201}.}
\label{fig:rest1}
\end{figure}

El reparto distinto entre ejecuciones no es un error. Las solicitudes que llegan mientras la primera espera la confirmación de Kafka, unos 4 ms, reciben \cod{409}, y las que llegan después reciben \cod{200} con la respuesta guardada. La clave reutilizada con otro pedido recibió \cod{409 CLAVE\_REUTILIZADA} y la solicitud sin clave \cod{400 FALTA\_IDEMPOTENCY\_KEY}.

\textbf{Parte 2.} Detuvimos Kafka con \cod{docker stop kafka}, enviamos dos veces el mismo pedido, volvimos a levantar Kafka y reintentamos. La Figura~\ref{fig:tiempos} compara los tiempos de respuesta de las dos partes en escala logarítmica.

\begin{figure}[H]
\centering
\begin{tikzpicture}
\begin{axis}[
    width=12.2cm, height=7cm,
    xbar, bar width=0.38cm,
    xmode=log, log basis x=10,
    xmin=0.5, xmax=40000,
    xtick={1,10,100,1000,10000}, xticklabels={1, 10, 100, 1000, 10000},
    xlabel={\small tiempo de respuesta en ms (escala logarítmica)},
    symbolic y coords={
        {503, Kafka caído},
        {201, reintento con Kafka de vuelta},
        {201, primer pedido},
        {201, pedidos siguientes},
        {409, otra igual en curso},
        {200 GET /menu, Kafka caído},
        {200, reintento},
    },
    ytick=data, y dir=reverse,
    tick label style={font=\scriptsize},
    nodes near coords, nodes near coords style={font=\scriptsize, anchor=west},
    point meta=explicit symbolic,
    enlarge y limits=0.13,
]
    \addplot[fill=rest!60, draw=rest] coordinates {
        (132,{201, primer pedido}) [132 ms]
        (5,{201, pedidos siguientes}) [5 ms]
        (1,{200, reintento}) [1 ms]
        (25,{201, reintento con Kafka de vuelta}) [25 ms]
        (1.3,{409, otra igual en curso}) [1,3 ms]
        (1.2,{200 GET /menu, Kafka caído}) [1,2 ms]
    };
    \addplot[fill=falla!55, draw=falla] coordinates {
        (5000,{503, Kafka caído}) [5000 ms]
    };
\end{axis}
\end{tikzpicture}
\caption{Tiempos de respuesta medidos. Con Kafka caído la API espera los 5 s de su límite antes de responder \cod{503}, pero las demás rutas siguen respondiendo en un milisegundo. Los \cod{201} con cocina, caja y facturación detenidas muestran que la API no espera a ningún servicio interno.}
\label{fig:tiempos}
\end{figure}

El intento fallido y el reintento usaron el mismo \cod{idPedido}, que empieza con \cod{52a1efde}, y en \cod{pedido.creado} quedó un solo evento. Esto importa porque un \cod{503} no prueba que Kafka no haya escrito el mensaje, solo que no llegó la confirmación. Si el reintento generara otro \cod{idPedido}, los consumidores verían dos pedidos distintos.

\subsection{Kafka: consumidor detenido}

Con Kafka y la API funcionando detuvimos Cocina, creamos un pedido desde Swagger y volvimos a levantar Cocina. La Figura~\ref{fig:kafka} muestra la línea de tiempo con las horas del log.

\begin{figure}[H]
\centering
\begin{tikzpicture}[x=0.075cm, y=1cm]
    % eje de 20:32:30 (0) a 20:35:30 (180 s)
    \fill[falla!10] (0,-0.35) rectangle (159,0.35);
    \node[font=\scriptsize, text=falla] at (60,0.55) {Cocina detenida};
    \fill[ws!15] (159,-0.35) rectangle (180,0.35);
    \node[font=\scriptsize, text=ws, anchor=south] at (170,0.35) {Cocina activa};
    \draw[flecha, gris] (0,-0.35) -- (185,-0.35);
    \foreach \t/\lab in {0/20:32:30, 60/20:33:30, 120/20:34:30, 180/20:35:30} {
        \draw[gris] (\t,-0.3) -- (\t,-0.4) node[below, font=\tiny]{\lab};
    }
    \fill[rest] (19,0) circle (3pt);
    \node[font=\scriptsize, text=rest, align=center, above] at (19,0.1) {\cod{201}\\pedido creado};
    \fill[kafka] (159,0) circle (3pt);
    \fill[kafka] (164,0) circle (3pt);
    \node[font=\scriptsize, text=kafka, align=left, anchor=north east] at (162,-0.75) {offset 8: \cod{EN\_PREPARACION}\\offset 9: \cod{LISTO}};
\end{tikzpicture}
\caption{Prueba de Kafka. El pedido \cod{3742ad17-...} se creó a las 20:32:49 con Cocina detenida. Al volver, Cocina publicó \cod{EN\_PREPARACION} a las 20:35:09 y \cod{LISTO} a las 20:35:14, los dos en la partición 0. El pedido esperó más de dos minutos en Kafka y no se perdió.}
\label{fig:kafka}
\end{figure}

El pedido no se perdió porque Kafka guarda para cada consumer group el offset hasta el que ya procesó, y Cocina hace el commit recién después de procesar. Al reiniciar continuó desde su último offset confirmado. Es el desacoplamiento temporal del Capítulo 2, donde productor y consumidor no necesitan estar activos al mismo tiempo \parencite{timbi2026comunicacion}.

\subsection{gRPC: Caja sin responder}

Domenica probó cuatro casos. En el primero llamó a un puerto sin servidor gRPC, y la llamada terminó en 2,02 s con \cod{DEADLINE\_EXCEEDED} en lugar de bloquear a Facturación (Figura~\ref{fig:deadline}).

\begin{figure}[H]
\centering
\begin{tikzpicture}
\begin{axis}[
    width=10cm, height=3.6cm,
    xbar, bar width=0.4cm,
    xmin=0, xmax=2.6,
    symbolic y coords={observado, deadline},
    ytick=data, tick label style={font=\small},
    xlabel={\small segundos},
    nodes near coords, nodes near coords style={font=\scriptsize, anchor=west},
    point meta=explicit symbolic,
    enlarge y limits=0.5,
]
    \addplot[fill=grpc!60, draw=grpc] coordinates {(2.00,deadline) [2,00 s] (2.02,observado) [2,02 s]};
\end{axis}
\end{tikzpicture}
\caption{Duración de la llamada gRPC sin servidor. Termina prácticamente en el deadline configurado.}
\label{fig:deadline}
\end{figure}

\begin{table}[H]
\centering
\small
\begin{tabularx}{\textwidth}{@{}c Y Y@{}}
\toprule
\textbf{\#} & \textbf{Caso} & \textbf{Resultado} \\
\midrule
1 & Llamada a un puerto sin servidor & \cod{DEADLINE\_EXCEEDED} en 2,02 s \\
2 & Pago confirmado y facturación & \cod{PAGADO} con \cod{PAG-3F40D3BBC3CA}, luego \cod{FACTURADO} con \cod{FAC-04CAAF4A013D} \\
3 & Pago y confirmación repetidos & La misma \cod{PAG-3F40D3BBC3CA}, \cod{repetida = true} y un solo \cod{PAGO/PAGADO} en Kafka \\
4 & Pago tardío, después de los 6 intentos & Con el diseño anterior quedó \cod{FACTURACION\_PENDIENTE}. Con el nuevo, Facturación leyó el \cod{PAGADO} guardado en Kafka y emitió \cod{FAC-6242528FFF56} \\
\bottomrule
\end{tabularx}
\caption{Casos de la prueba de gRPC. El caso 4 fue un resultado inesperado que nos llevó a la versión 1.3 del contrato (Figura~\ref{fig:v13}).}
\label{tab:grpc}
\end{table}

\subsection{WebSocket: cliente desconectado}

Cinthya probó que un estudiante que pierde la conexión no pierde el estado de su pedido. El cliente creó un pedido, se desconectó antes de que Cocina terminara y volvió quince segundos después.

\begin{figure}[H]
\centering
\begin{tikzpicture}[x=0.52cm, y=1cm]
    \fill[ws!15] (0,-0.3) rectangle (2.5,0.3);
    \fill[falla!10] (2.5,-0.3) rectangle (17.5,0.3);
    \fill[ws!15] (17.5,-0.3) rectangle (27,0.3);
    \node[font=\scriptsize, text=ws] at (1.25,0) {conectado};
    \node[font=\scriptsize, text=falla] at (10,0) {desconectado 15 s};
    \node[font=\scriptsize, text=ws] at (23.5,0) {conectado};
    \draw[flecha, gris] (0,-0.3) -- (27.8,-0.3) node[right, font=\scriptsize]{s};
    \foreach \t in {0,5,10,15,20,25} {\draw[gris] (\t,-0.25) -- (\t,-0.35) node[below, font=\tiny]{\t};}

    \foreach \t/\y/\txt/\col in {
        0.3/2.15/{\cod{POST} 201 y WebSocket abierto}/rest,
        0.5/1.45/{suscribir, luego \cod{GET}: \cod{EN\_PREPARACION}}/ws,
        2.5/0.75/{el cliente se desconecta}/falla,
        17.5/2.15/{vuelve y se suscribe}/ws,
        17.7/1.45/{\cod{GET}: \cod{LISTO}}/rest,
        20.5/0.75/{push \cod{ENTREGADO}}/ws} {
        \draw[\col] (\t,0.3) -- (\t,\y);
        \node[font=\scriptsize, text=\col, anchor=west, inner sep=1.5pt] at (\t,\y) {\txt};
    }
    \node[font=\scriptsize, text=kafka, align=center] at (10,-1.05) {mientras tanto Cocina publica \cod{LISTO}\\y Notificaciones descarta el aviso};
\end{tikzpicture}
\caption{Línea de tiempo real de la prueba de WebSocket. El estado perdido durante la desconexión se recuperó con \cod{GET} y el evento siguiente llegó por push.}
\label{fig:ws}
\end{figure}

El WebSocket es comunicación transitoria, así que el aviso de \cod{LISTO} se descartó porque nadie estaba conectado. Lo que no se perdió fue el estado, porque Kafka es persistente, la API consumió el evento y lo devolvió por REST. Combinar las dos vías da inmediatez cuando hay conexión y consistencia cuando no la hay.

\subsection{Integración: flujo completo}

Con Kafka y los seis servicios levantados, un script hizo lo mismo que el estudiante y el cajero. Creó un pedido de 450 centavos, se suscribió, intentó entregarlo antes de tiempo, lo pagó dos veces, intentó pagar un pedido inexistente y lo entregó dos veces. La Figura~\ref{fig:integracion} ubica cada paso en el tiempo y en su mecanismo.

\begin{figure}[H]
\centering
\begin{tikzpicture}[
    xscale=0.95,
    celda/.style={font=\scriptsize, align=center, inner sep=2pt},
    pto/.style={circle, inner sep=0pt, minimum size=5pt}
]
    \foreach \y/\m/\col in {3/REST/rest, 2/Kafka/kafka, 1/gRPC/grpc, 0/WebSocket/ws} {
        \node[font=\small\bfseries, text=\col, anchor=east] at (-0.4,\y*1.15) {\m};
        \draw[\col!30] (-0.2,\y*1.15) -- (13.4,\y*1.15);
    }
    \foreach \x/\t in {0.8/{+0,1 s}, 5.0/{+5,1 s}, 7.6/{+5,3 s}, 10.0/{+5,4 s}, 12.4/{+5,5 s}} {
        \node[font=\scriptsize\bfseries, text=gris] at (\x,4.35) {\t};
        \draw[gris!40, dashed] (\x,4.1) -- (\x,-0.4);
    }
    \draw[decorate, gris] (1.8,4.35) -- (4.0,4.35);
    \node[font=\scriptsize, text=gris, fill=white] at (2.9,4.35) {Cocina 5 s};

    % REST
    \node[celda, text=rest, fill=white] at (0.8,3.45) {\cod{201} creado\\\cod{409} no listo};
    \node[celda, text=falla, fill=white] at (7.6,3.45) {\cod{409}\\no pagado};
    \node[celda, text=rest, fill=white] at (10.0,3.45) {pago dos veces\\misma confirmación};
    \node[celda, text=rest, fill=white] at (12.4,3.45) {entrega dos\\veces, \cod{200}};
    % Kafka
    \node[celda, text=kafka, fill=white] at (0.8,2.3) {\cod{EN\_PREPARACION}\\\cod{PENDIENTE}};
    \node[celda, text=kafka, fill=white] at (5.0,2.3) {\cod{LISTO}};
    \node[celda, text=kafka, fill=white] at (10.0,2.3) {\cod{PAGADO}\\\cod{FACTURADO}};
    \node[celda, text=kafka, fill=white] at (12.4,2.3) {\cod{ENTREGADO}};
    % gRPC
    \node[celda, text=grpc, fill=white] at (10.0,1.15) {\cod{ConfirmarPago}};
    % WebSocket
    \node[celda, text=ws, fill=white] at (0.8,0) {suscrito};
    \node[celda, text=ws, fill=white] at (5.0,0) {push \cod{LISTO}};
    \node[celda, text=ws, fill=white] at (12.4,0) {6 estados\\recibidos};
\end{tikzpicture}
\caption{Línea de tiempo de la prueba de integración, por mecanismo. La API respondió en 0,1 s mientras Cocina tardaba 5 s, y en los 0,4 s siguientes a \cod{LISTO} se completaron el pago, la confirmación por gRPC, la factura y la entrega.}
\label{fig:integracion}
\end{figure}

\begin{table}[H]
\centering
\small
\begin{tabular}{@{}l c l l@{}}
\toprule
\textbf{Topic} & \textbf{Offset} & \textbf{Evento} & \textbf{Publicado por} \\
\midrule
\cod{pedido.creado} & 7 & pedido creado, total 450 & API \\
\cod{pedido.estado} & 27 & \cod{PREPARACION/EN\_PREPARACION} & Cocina \\
\cod{pedido.estado} & 28 & \cod{PAGO/PENDIENTE} & Caja \\
\cod{pedido.estado} & 29 & \cod{PREPARACION/LISTO} & Cocina \\
\cod{pedido.estado} & 30 & \cod{PAGO/PAGADO} & Caja \\
\cod{pedido.estado} & 31 & \cod{FACTURACION/FACTURADO} & Facturación \\
\cod{pedido.estado} & 32 & \cod{ENTREGA/ENTREGADO} & API \\
\bottomrule
\end{tabular}
\caption{Eventos del pedido \cod{0751c38d-...} en Kafka. Los siete quedaron en la partición 0 y en orden, aunque los publicaron cuatro servicios distintos. Pese a pagar y entregar dos veces, hay un solo \cod{PAGADO} y un solo \cod{ENTREGADO}.}
\label{tab:offsets}
\end{table}

\begin{figure}[H]
\centering
\begin{tikzpicture}
\begin{axis}[
    width=7.6cm, height=5.4cm,
    ybar, bar width=0.45cm,
    ymin=0, ymax=23,
    symbolic x coords={Recibidos, Pagados, Entregados},
    xtick=data, enlarge x limits=0.25,
    nodes near coords, nodes near coords style={font=\scriptsize},
    tick label style={font=\small},
    legend style={font=\scriptsize, at={(0.98,0.97)}, anchor=north east, draw=none},
    ylabel={\small pedidos}, title={\small Contadores de Estadísticas},
]
    \addplot[fill=gris!30, draw=gris] coordinates {(Recibidos,18) (Pagados,3) (Entregados,3)};
    \addplot[fill=kafka!55, draw=kafka] coordinates {(Recibidos,19) (Pagados,4) (Entregados,4)};
    \legend{antes, después}
\end{axis}
\end{tikzpicture}
\hfill
\begin{tikzpicture}
\begin{axis}[
    width=5.4cm, height=5.4cm,
    ybar, bar width=0.7cm,
    ymin=0, ymax=2100,
    symbolic x coords={antes, después},
    xtick=data, enlarge x limits=0.5,
    nodes near coords, nodes near coords style={font=\scriptsize},
    tick label style={font=\small},
    ylabel={\small centavos}, title={\small Ventas entregadas},
]
    \addplot[fill=kafka!55, draw=kafka] coordinates {(antes,1350) (después,1800)};
\end{axis}
\end{tikzpicture}
\caption{Estadísticas antes y después del pedido de la prueba. Subió en uno cada contador y las ventas aumentaron justo los 450 centavos del pedido, aunque el pago y la entrega se enviaron dos veces.}
\label{fig:estadisticas}
\end{figure}

La prueba relaciona varios conceptos del Capítulo 2. REST es síncrono porque el estudiante y el cajero necesitan una respuesta inmediata, mientras que Kafka desacopla a los servicios en el tiempo y en la identidad, porque el productor no sabe quién consume. gRPC se usó en el único punto donde hacía falta una respuesta confirmada, y el WebSocket empujó cada cambio al estudiante. El flujo funcionó sin ajustes porque los cuatro integrantes respetamos el mismo contrato.

\section{Conclusiones}

No hay un mecanismo que sirva para todo, y el proyecto lo mostró en un caso concreto. REST funcionó bien donde una persona espera una respuesta, Kafka donde los servicios deben seguir trabajando aunque otro se caiga, gRPC en la única consulta interna que necesitaba una confirmación inmediata y WebSocket para avisar al estudiante sin que tenga que preguntar.

Las fallas que pedía la consigna se resolvieron con decisiones de diseño más que con código extra. La clave de idempotencia evitó pedidos duplicados, esperar la confirmación de Kafka evitó responder \cod{201} sin evento, el deadline evitó que Facturación se bloqueara y el orden de suscribirse antes de consultar evitó que el estudiante perdiera un cambio. En todos los casos combinar la semántica al menos una vez del cliente con la de a lo sumo una vez del servidor dejó un solo efecto.

El resultado inesperado de la prueba de gRPC fue lo que más nos enseñó. Usar una llamada síncrona para esperar algo que ocurre cuando el estudiante quiere dejaba pedidos sin factura, y la solución fue combinar los dos tipos de comunicación, un evento asíncrono que avisa y una llamada síncrona que confirma.

Las principales limitaciones son que los datos viven en memoria, que la API funciona como una sola instancia y que no hay autenticación. Kafka también es un punto único, y con el broker caído no se pueden crear pedidos.

\subsection{Uso de inteligencia artificial}

Usamos herramientas de IA como apoyo. Las consultas y la reflexión de cada integrante están en \cod{docs/bitacora-ia.md}, y los archivos de código hechos con ayuda de IA llevan una marca con el nombre de quien los revisó.

\subsection{Cómo ejecutar el sistema}

Todos los comandos se ejecutan desde la raíz del repositorio, con el entorno virtual activado y cada servicio en su propia terminal.

\begin{table}[H]
\centering
\small
\begin{tabularx}{\textwidth}{@{}l Y@{}}
\toprule
\textbf{Servicio} & \textbf{Comando} \\
\midrule
Kafka & \cod{docker compose -f infra/docker-compose.yml up -d} \\
API de Pedidos & \cod{python -m backend.api\_rest.main} \\
Cocina & \cod{python -m backend.cocina.main --demora 5} \\
Estadísticas & \cod{python -m backend.estadisticas.main} \\
Caja & \cod{python -m uvicorn backend.caja.main:app --port 8001} \\
Facturación & \cod{python -m backend.facturacion.main} \\
Notificaciones & \cod{python -m backend.websockets.servidor} \\
Cliente web & abrir \cod{backend/websockets/cliente/index.html} \\
Pruebas automáticas & \cod{python -m pytest -q}, 62 pruebas \\
Prueba de integración & \cod{python -m backend.api\_rest.pruebas.prueba\_flujo\_completo} \\
\bottomrule
\end{tabularx}
\caption{Comandos para levantar y probar el sistema.}
\label{tab:ejecucion}
\end{table}

El código, los ADR completos, las pruebas con sus evidencias y la bitácora están en el repositorio del grupo, \href{https://github.com/DavidLarriva/Mecanismo-comunicacion-P02}{\cod{github.com/DavidLarriva/Mecanismo-comunicacion-P02}}.



% ==========================================
% REFERENCIAS BIBLIOGRÁFICAS
% ==========================================

\printbibliography[title={Referencias}]

\end{document}
